% English interface strings
\newcommand\ACNstrCourse{Advanced Communication Networks}
\newcommand\ACNstrNotes{Lecture Notes}
\newcommand\ACNstrLecture{Lecture}
\newcommand\ACNstrContents{In this lecture}
\newcommand\ACNstrPart{Course part}
\newcommand\ACNstrTime{Estimated study time}
\newcommand\ACNstrLevel{Difficulty}
\newcommand\ACNstrPrereq{Prerequisites}
\newcommand\ACNstrDefinition{Definition}
\newcommand\ACNstrResult{Result}
\newcommand\ACNstrExample{Worked example}
\newcommand\ACNstrNote{Engineering note}
\newcommand\ACNstrWarning{Common pitfall}
\newcommand\ACNstrTakeaways{Key takeaways}
\newcommand\ACNstrProcedure{Procedure}
\newcommand\ACNstrExercise{Exercise}
\newcommand\ACNstrSolution{Solution to Exercise}
\newcommand\ACNstrLevelOne{Level 1 -- conceptual}
\newcommand\ACNstrLevelTwo{Level 2 -- numerical}
\newcommand\ACNstrLevelThree{Level 3 -- analytical}
\newcommand\ACNstrKTone{Term}
\newcommand\ACNstrKTtwo{Persian}
\newcommand\ACNstrKTthree{Meaning in this lecture}
\newcommand\ACNstrReview{To be verified}
\newcommand\ACNstrCoverNote{Original lecture notes for the course \emph{Advanced
  Communication Networks}. The lecture sequence follows the course syllabus
  (26 lectures, 2G to 6G). Equations use identical notation in the English and
  Persian editions. Companion edition: Persian (\texttt{lecture-\ACNnumber-fa.pdf}).}
\newcommand\ACNstrRefs{References}
\crefname{equation}{Eq.}{Eqs.}
\Crefname{equation}{Equation}{Equations}
\crefname{figure}{Figure}{Figures}
\crefname{table}{Table}{Tables}
\crefname{section}{Section}{Sections}
\crefname{subsection}{Section}{Sections}
\crefname{acnexample}{Example}{Examples}
\crefname{acnexercise}{Exercise}{Exercises}
\crefname{acndefinition}{Definition}{Definitions}
\crefname{acntheorem}{Result}{Results}
\crefname{acnprocedure}{Procedure}{Procedures}
